% ============================================================
%  TEMA 2 — APARTADO 1: ANÁLISIS DE LOS PUESTOS DE TRABAJO
%
%  RECORDATORIO AL EQUIPO:
%  Las etapas del análisis de puestos (ADP) son generalmente:
%    1. Planificación del ADP (objetivos, recursos)
%    2. Preparación (diseño de instrumentos: cuestionarios, entrevistas...)
%    3. Recogida de información (observación, entrevista, cuestionario, diario...)
%    4. Análisis y tratamiento de datos
%    5. Elaboración de la descripción y especificación del puesto
%    6. Actualización y mantenimiento
%
%  PARA CADA ETAPA INDICAD:
%    - ¿La realiza CaixaBank? ¿Cómo exactamente?
%    - Ventajas e inconvenientes de su enfoque
%    - ¿Está la empresa satisfecha con ese proceso?
%    - ¿Habría maneras más eficientes de hacerlo?
% ============================================================

\section{Tema 2: Análisis de los puestos de trabajo}

\subsection{Etapas del análisis de puestos desarrolladas en CaixaBank}

El análisis de puestos de trabajo (APT) es el proceso sistemático de recopilar,
analizar y organizar información sobre los puestos de una organización con el
objetivo de elaborar descripciones y especificaciones de puestos que sirvan de
base para múltiples decisiones de RRHH (selección, formación, retribución,
evaluación del desempeño, etc.).

A continuación se detalla en qué medida CaixaBank desarrolla cada una de las
etapas clásicas del APT:

% ---- TABLA RESUMEN DE ETAPAS ----
\begin{table}[H]
  \centering
  \small
  \caption{Resumen de etapas del APT en CaixaBank}
  \begin{tabularx}{\textwidth}{p{0.28\textwidth} p{0.1\textwidth} X}
    \toprule
    \textbf{Etapa} & \textbf{¿Se realiza?} & \textbf{Observaciones} \\
    \midrule
    1. Planificación del APT    & [Sí/No/Parcial] & [Completar con información de la entrevista] \\
    2. Preparación de instrumentos & [Sí/No/Parcial] & \\
    3. Recogida de información  & [Sí/No/Parcial] & \\
    4. Análisis y tratamiento   & [Sí/No/Parcial] & \\
    5. Descripción y especificación & [Sí/No/Parcial] & \\
    6. Actualización y mantenimiento & [Sí/No/Parcial] & \\
    \bottomrule
  \end{tabularx}
\end{table}

% ---- ETAPA 1 ----
\subsubsection{Etapa 1 — Planificación del análisis de puestos}

% TODO: Describir si CaixaBank planifica formalmente el APT.
% ¿Tiene un departamento o equipo dedicado? ¿Con qué periodicidad lo actualiza?
% ¿Lo hace internamente o externaliza (consultoría)?

[Descripción a completar tras la entrevista]

\textbf{Ventajas del enfoque adoptado:}
\begin{itemize}
  \item [Completar]
\end{itemize}

\textbf{Inconvenientes / áreas de mejora:}
\begin{itemize}
  \item [Completar]
\end{itemize}

% ---- ETAPA 2 ----
\subsubsection{Etapa 2 — Preparación de los instrumentos de recogida}

% TODO: ¿Qué métodos de recogida de información utiliza CaixaBank?
% Métodos habituales: cuestionario estructurado, entrevista al ocupante,
% entrevista al supervisor, observación directa, diario de actividades,
% incidentes críticos, conferencia técnica de expertos.
% Indicad cuáles SÍ usa, cuáles NO y por qué (coste, tipo de puesto, etc.)

[Descripción a completar]

% ---- ETAPA 3 ----
\subsubsection{Etapa 3 — Recogida de información}

% TODO: ¿Quién recoge la información? ¿El propio empleado, su supervisor,
% un analista de RRHH? ¿Hay sesgo? ¿Cómo garantizan la fiabilidad?

[Descripción a completar]

% ---- ETAPAS 4, 5 y 6 ----
\subsubsection{Etapas 4, 5 y 6 — Análisis, descripción y actualización}

% TODO: ¿Existe un catálogo formal de puestos? ¿Con qué frecuencia se revisa?
% En un banco como CaixaBank la aparición de nuevas figuras (analista de datos,
% ciberseguridad, UX bancario) exige actualizaciones frecuentes.
% ¿Lo hace CaixaBank? ¿Cómo gestionan la obsolescencia de los perfiles?

[Descripción a completar]

\begin{notaequipo}
  Preguntad explícitamente en la entrevista:
  \begin{enumerate}
    \item ¿Tienen un manual o catálogo de puestos formalizado?
    \item ¿Quién es responsable de actualizar las descripciones de puesto?
    \item ¿Con qué frecuencia se revisan? ¿Se revisaron tras la fusión con Bankia?
    \item ¿Cómo recogen la información del puesto (cuestionarios, observación...)?
    \item ¿Están satisfechos con el proceso actual? ¿Qué mejorarían?
  \end{enumerate}
\end{notaequipo}
